% =============================================================================
% Chapter 14: SQL - Λυμένες Ασκήσεις
% Data Mining Study Guide - NTUA
% =============================================================================

\chapter{SQL - Λυμένες Ασκήσεις}
\label{ch:sql-exercises}

Αυτό το κεφάλαιο περιέχει λυμένες ασκήσεις για SQL queries.

% -----------------------------------------------------------------------------
% Exercise 3.1
% -----------------------------------------------------------------------------
\section{Άσκηση 3.1: SELECT με WHERE}

\begin{formulabox}[Εκφώνηση]
Δίνεται ο πίνακας \texttt{Products(pid, pname, category, price, stock)}.

Γράψτε SQL query για να βρείτε τα προϊόντα κατηγορίας «Electronics» με τιμή μεταξύ
100 και 500 ευρώ, ταξινομημένα κατά τιμή (φθίνουσα).
\end{formulabox}

\subsection*{Λύση}

\begin{lstlisting}[language=SQL]
SELECT pid, pname, price
FROM Products
WHERE category = 'Electronics'
  AND price BETWEEN 100 AND 500
ORDER BY price DESC;
\end{lstlisting}

\begin{infobox}[BETWEEN]
Το \texttt{BETWEEN 100 AND 500} είναι ισοδύναμο με \texttt{price >= 100 AND price <= 500}
(περιλαμβάνει τα όρια).
\end{infobox}

% -----------------------------------------------------------------------------
% Exercise 3.2
% -----------------------------------------------------------------------------
\section{Άσκηση 3.2: JOINs (INNER, LEFT)}

\begin{formulabox}[Εκφώνηση]
Δίνονται:
\begin{itemize}
    \item \texttt{Customers(cid, cname, city)}
    \item \texttt{Orders(oid, cid, order\_date, total)}
\end{itemize}

\begin{enumerate}
    \item Βρείτε όλους τους πελάτες με τις παραγγελίες τους
    \item Βρείτε όλους τους πελάτες, ακόμα κι αν δεν έχουν παραγγελίες
\end{enumerate}
\end{formulabox}

\subsection*{Λύση}

\textbf{(a) INNER JOIN -- μόνο πελάτες με παραγγελίες:}

\begin{lstlisting}[language=SQL]
SELECT c.cname, o.oid, o.total
FROM Customers c
INNER JOIN Orders o ON c.cid = o.cid;
\end{lstlisting}

\textbf{(b) LEFT JOIN -- όλοι οι πελάτες:}

\begin{lstlisting}[language=SQL]
SELECT c.cname, o.oid, o.total
FROM Customers c
LEFT JOIN Orders o ON c.cid = o.cid;
\end{lstlisting}

\begin{center}
\begin{tikzpicture}[scale=0.8]
    % INNER JOIN
    \begin{scope}[xshift=-4cm]
        \draw[fill=box-blue, fill opacity=0.5] (0,0) circle (1.2);
        \draw[fill=box-green, fill opacity=0.5] (1.5,0) circle (1.2);
        \fill[secondary, opacity=0.7] (0.75,0) circle (0.5);
        \node at (0.75,-2) {\small INNER JOIN};
    \end{scope}
    % LEFT JOIN
    \begin{scope}[xshift=4cm]
        \draw[fill=secondary, fill opacity=0.5] (0,0) circle (1.2);
        \draw[fill=box-green, fill opacity=0.3] (1.5,0) circle (1.2);
        \node at (0.75,-2) {\small LEFT JOIN};
    \end{scope}
\end{tikzpicture}
\end{center}

% -----------------------------------------------------------------------------
% Exercise 3.3
% -----------------------------------------------------------------------------
\section{Άσκηση 3.3: GROUP BY με COUNT}

\begin{formulabox}[Εκφώνηση]
Χρησιμοποιώντας τους πίνακες \texttt{Orders} και \texttt{Customers}, βρείτε πόσες
παραγγελίες έχει κάνει κάθε πελάτης.
\end{formulabox}

\subsection*{Λύση}

\begin{lstlisting}[language=SQL]
SELECT c.cid, c.cname, COUNT(o.oid) AS order_count
FROM Customers c
LEFT JOIN Orders o ON c.cid = o.cid
GROUP BY c.cid, c.cname;
\end{lstlisting}

\begin{warningbox}[LEFT vs INNER JOIN]
Χρησιμοποιούμε LEFT JOIN για να συμπεριλάβουμε πελάτες με 0 παραγγελίες.
Με INNER JOIN θα χάναμε αυτούς τους πελάτες!
\end{warningbox}

% -----------------------------------------------------------------------------
% Exercise 3.4
% -----------------------------------------------------------------------------
\section{Άσκηση 3.4: GROUP BY με HAVING}

\begin{formulabox}[Εκφώνηση]
Βρείτε τις πόλεις που έχουν περισσότερους από 5 πελάτες και εμφανίστε τον αριθμό
πελατών ανά πόλη.
\end{formulabox}

\subsection*{Λύση}

\begin{lstlisting}[language=SQL]
SELECT city, COUNT(*) AS customer_count
FROM Customers
GROUP BY city
HAVING COUNT(*) > 5
ORDER BY customer_count DESC;
\end{lstlisting}

\begin{infobox}[WHERE vs HAVING]
\begin{itemize}
    \item \textbf{WHERE}: Φιλτράρει \textit{γραμμές} \textbf{πριν} το grouping
    \item \textbf{HAVING}: Φιλτράρει \textit{groups} \textbf{μετά} το grouping
\end{itemize}
\end{infobox}

% -----------------------------------------------------------------------------
% Exercise 3.5
% -----------------------------------------------------------------------------
\section{Άσκηση 3.5: Subquery με IN}

\begin{formulabox}[Εκφώνηση]
Βρείτε τους πελάτες που έχουν κάνει παραγγελία με total > 1000.
\end{formulabox}

\subsection*{Λύση}

\textbf{Μέθοδος 1: Subquery με IN}

\begin{lstlisting}[language=SQL]
SELECT cname
FROM Customers
WHERE cid IN (
    SELECT cid
    FROM Orders
    WHERE total > 1000
);
\end{lstlisting}

\textbf{Μέθοδος 2: JOIN (ισοδύναμο)}

\begin{lstlisting}[language=SQL]
SELECT DISTINCT c.cname
FROM Customers c
JOIN Orders o ON c.cid = o.cid
WHERE o.total > 1000;
\end{lstlisting}

% -----------------------------------------------------------------------------
% Exercise 3.6
% -----------------------------------------------------------------------------
\section{Άσκηση 3.6: Correlated Subquery}

\begin{formulabox}[Εκφώνηση]
Βρείτε τους υπαλλήλους που έχουν μισθό μεγαλύτερο από τον μέσο όρο του τμήματός τους.
\end{formulabox}

\subsection*{Λύση}

\begin{lstlisting}[language=SQL]
SELECT e1.ename, e1.salary, e1.dept_id
FROM Employee e1
WHERE e1.salary > (
    SELECT AVG(e2.salary)
    FROM Employee e2
    WHERE e2.dept_id = e1.dept_id  -- Correlated!
);
\end{lstlisting}

\begin{infobox}[Correlated Subquery]
Το subquery εξαρτάται από το εξωτερικό query (αναφέρεται στο \texttt{e1.dept\_id}).
Εκτελείται μία φορά \textbf{για κάθε γραμμή} του εξωτερικού query.
\end{infobox}

% -----------------------------------------------------------------------------
% Exercise 3.7
% -----------------------------------------------------------------------------
\section{Άσκηση 3.7: NOT EXISTS Pattern}

\begin{formulabox}[Εκφώνηση]
Βρείτε τους πελάτες που \textbf{δεν} έχουν κάνει καμία παραγγελία.
\end{formulabox}

\subsection*{Λύση}

\textbf{Μέθοδος 1: NOT EXISTS}

\begin{lstlisting}[language=SQL]
SELECT c.cname
FROM Customers c
WHERE NOT EXISTS (
    SELECT 1
    FROM Orders o
    WHERE o.cid = c.cid
);
\end{lstlisting}

\textbf{Μέθοδος 2: LEFT JOIN + IS NULL}

\begin{lstlisting}[language=SQL]
SELECT c.cname
FROM Customers c
LEFT JOIN Orders o ON c.cid = o.cid
WHERE o.oid IS NULL;
\end{lstlisting}

\textbf{Μέθοδος 3: NOT IN}

\begin{lstlisting}[language=SQL]
SELECT cname
FROM Customers
WHERE cid NOT IN (SELECT cid FROM Orders);
\end{lstlisting}

\begin{warningbox}[NOT IN με NULL]
Αν το subquery του NOT IN επιστρέψει NULL, το αποτέλεσμα είναι κενό!
Προτιμήστε NOT EXISTS για ασφάλεια.
\end{warningbox}

% -----------------------------------------------------------------------------
% Exercise 3.8
% -----------------------------------------------------------------------------
\section{Άσκηση 3.8: SQL Division (Double NOT EXISTS)}

\begin{formulabox}[Εκφώνηση]
Δίνονται:
\begin{itemize}
    \item \texttt{Enrollment(sid, cid)}
    \item \texttt{Course(cid, cname)}
\end{itemize}

Βρείτε τους φοιτητές που έχουν εγγραφεί σε \textbf{όλα} τα μαθήματα.
\end{formulabox}

\subsection*{Λύση}

\textbf{Λογική:} «Δεν υπάρχει μάθημα στο οποίο ο φοιτητής να μην έχει εγγραφεί»

\begin{lstlisting}[language=SQL]
SELECT DISTINCT e1.sid
FROM Enrollment e1
WHERE NOT EXISTS (
    -- Μαθήματα που ΔΕΝ έχει πάρει ο e1.sid
    SELECT c.cid
    FROM Course c
    WHERE NOT EXISTS (
        -- Έλεγχος αν ο e1.sid έχει πάρει το c.cid
        SELECT 1
        FROM Enrollment e2
        WHERE e2.sid = e1.sid AND e2.cid = c.cid
    )
);
\end{lstlisting}

\begin{infobox}[Ερμηνεία Double NOT EXISTS]
\begin{enumerate}
    \item Εσωτερικό NOT EXISTS: Βρες μαθήματα που \textbf{δεν} έχει πάρει ο φοιτητής
    \item Εξωτερικό NOT EXISTS: Κράτα φοιτητές που \textbf{δεν} έχουν τέτοια μαθήματα
\end{enumerate}
Δηλαδή: φοιτητές χωρίς «κενά» στο πρόγραμμα.
\end{infobox}

% -----------------------------------------------------------------------------
% Exercise 3.9
% -----------------------------------------------------------------------------
\section{Άσκηση 3.9: Multiple JOINs με Aggregation}

\begin{formulabox}[Εκφώνηση]
Δίνονται:
\begin{itemize}
    \item \texttt{Orders(oid, cid, order\_date)}
    \item \texttt{OrderItems(oid, pid, quantity, unit\_price)}
    \item \texttt{Products(pid, pname, category)}
\end{itemize}

Βρείτε το συνολικό έσοδο (revenue) ανά κατηγορία προϊόντος.
\end{formulabox}

\subsection*{Λύση}

\begin{lstlisting}[language=SQL]
SELECT
    p.category,
    SUM(oi.quantity * oi.unit_price) AS total_revenue
FROM Products p
JOIN OrderItems oi ON p.pid = oi.pid
GROUP BY p.category
ORDER BY total_revenue DESC;
\end{lstlisting}

% -----------------------------------------------------------------------------
% Exercise 3.10
% -----------------------------------------------------------------------------
\section{Άσκηση 3.10: Complex Query}

\begin{formulabox}[Εκφώνηση]
Βρείτε τους top 3 πελάτες (με βάση συνολικές αγορές) που έχουν αγοράσει προϊόντα
από τουλάχιστον 2 διαφορετικές κατηγορίες.
\end{formulabox}

\subsection*{Λύση}

\begin{lstlisting}[language=SQL]
SELECT
    c.cid,
    c.cname,
    SUM(oi.quantity * oi.unit_price) AS total_spent,
    COUNT(DISTINCT p.category) AS categories_bought
FROM Customers c
JOIN Orders o ON c.cid = o.cid
JOIN OrderItems oi ON o.oid = oi.oid
JOIN Products p ON oi.pid = p.pid
GROUP BY c.cid, c.cname
HAVING COUNT(DISTINCT p.category) >= 2
ORDER BY total_spent DESC
LIMIT 3;
\end{lstlisting}

\begin{infobox}[Σειρά Εκτέλεσης SQL]
\begin{enumerate}
    \item FROM / JOIN
    \item WHERE
    \item GROUP BY
    \item HAVING
    \item SELECT
    \item ORDER BY
    \item LIMIT
\end{enumerate}
\end{infobox}
